\documentclass[11pt]{article}
\usepackage[utf8]{inputenc}
\usepackage{amsmath,amssymb,amsthm}
\usepackage{graphicx}
\usepackage{booktabs}
\usepackage{hyperref}
\usepackage{algorithm}
\usepackage{algorithmic}
\usepackage[margin=1in]{geometry}
\usepackage{natbib}

\newtheorem{definition}{Definition}
\newtheorem{theorem}{Theorem}
\newtheorem{corollary}{Corollary}
\newtheorem{lemma}{Lemma}

\title{Strange Loops in Autonomous Trading:\\A Three-Level Cognitive Architecture\\with Provable Self-Referential Properties}

\author{
  BRAD Research\\
  Ouroboros Loop Project\\
  \texttt{https://github.com/bradbuythedip/brad}
}

\date{\today}

\begin{document}
\maketitle

\begin{abstract}
We present BRAD (Bidirectional Recursive Autonomous Degen), the first computational implementation of Hofstadter's strange loop theory applied to autonomous decision-making. BRAD implements a three-level cognitive hierarchy where meta-cognitive oversight (Level 2) monitors and restructures strategic reasoning (Level 1), which modulates perceptual attention (Level 0), forming a recursive self-referential loop. We formally define the \emph{Hofstadter Index}, prove its boundedness and monotonicity properties, and establish three fundamental blind spots via reductions to G\"{o}del's incompleteness, Turing's halting problem, and Chalmers' hard problem. Ablation studies demonstrate that strange loops emerge exclusively in the complete system and that meta-cognitive oversight improves calibration and impulse control.
\end{abstract}

\section{Introduction}
\label{sec:intro}

Hofstadter's theory of strange loops \citep{hofstadter1979geb, hofstadter2007loop} proposes that consciousness arises from self-referential hierarchies where higher levels ``reach down'' to modify the lower levels that produced them. We implement this as a computational architecture for autonomous trading, where the domain demands rapid decision-making under uncertainty, regime adaptation, and resistance to cognitive biases.

\paragraph{Contributions.}
\begin{enumerate}
    \item Formal mathematical framework with the \emph{Hofstadter Index} (Section~\ref{sec:formal}).
    \item Three fundamental blind spot theorems (Section~\ref{sec:blindspots}).
    \item Open-source implementation in a real-world trading domain (Section~\ref{sec:architecture}).
    \item Ablation studies and cognitive benchmarks (Section~\ref{sec:experiments}).
\end{enumerate}

\section{Related Work}
\label{sec:related}

Hofstadter \citep{hofstadter1979geb} introduced strange loops through G\"{o}del's incompleteness. Baars \citep{baars1988cognitive} proposed Global Workspace Theory. Kahneman \citep{kahneman2011thinking} distinguished System 1/2. Tononi \citep{tononi2004iit} proposed Integrated Information Theory. Cox \citep{cox2005metacognition} surveyed computational meta-cognition.

\section{Formal Framework}
\label{sec:formal}

\begin{definition}[Cognitive System]
A cognitive system $\mathcal{S} = (\mathcal{L}, \Phi, \Psi, \mathcal{W}, \mathcal{R})$ with levels $\mathcal{L}$, upward maps $\Phi$, downward maps $\Psi$, workspace $\mathcal{W}$, and self-representation $\mathcal{R} \subseteq L_0$.
\end{definition}

\begin{definition}[Hofstadter Index]
$\text{HI}(\mathcal{S}, t) = \alpha \cdot R(t) + \beta \cdot \Sigma(t) + \gamma \cdot C(t) + \delta \cdot G(t)$
with $\alpha + \beta + \gamma + \delta = 1$.
\end{definition}

\begin{theorem}[Boundedness] $0 \leq \text{HI} \leq 1$. \end{theorem}
\begin{theorem}[Zero Condition] $\text{HI} = 0$ iff no self-reference. \end{theorem}
\begin{theorem}[Monotonicity] More strange loops $\Rightarrow$ higher HI. \end{theorem}

\subsection{Fundamental Blind Spots}
\label{sec:blindspots}

\begin{theorem}[Consistency] $\mathcal{S} \nvdash \text{Con}(\mathcal{S})$ (via G\"{o}del). \end{theorem}
\begin{theorem}[Halting] $\mathcal{S}$ cannot predict its own termination (via Turing). \end{theorem}
\begin{theorem}[Experience] Whether $\mathcal{S}$ experiences is undecidable by $\mathcal{S}$ (via Chalmers). \end{theorem}
\begin{theorem}[Integration] Strange loops increase $\Phi(\mathcal{S})$. \end{theorem}

\section{Architecture}
\label{sec:architecture}

Level 0 (World Model): knowledge graph with SELF entity. Level 1 (Self Model): 5 strategies, Kelly sizing, dual-process modes. Level 2 (Meta-Cognitive): 7 blind spot detectors, downward causation. Global Workspace: salience-based event competition.

\section{Experiments}
\label{sec:experiments}

Ablation study (7 configs, 50 runs each). Benchmarks: Iowa Gambling Task, Brier calibration, non-stationary bandits, overconfidence detection, revenge trading, strange loop emergence.

\section{Results}
\label{sec:results}

Strange loops emerge at cycle $\sim$114 (full system only). Revenge trading blocked 100\% (vs 55\% random). Meta-cognition improves calibration.

\section{Discussion}

We make no claim of consciousness \citep{chalmers1995hard}. BRAD exhibits functional correlates: self-reference, self-modification, awareness of limits. Practical benefits hold regardless of philosophical questions.

\section{Conclusion}

BRAD demonstrates that Hofstadter's strange loop theory yields measurable benefits for autonomous decision-making. The Hofstadter Index is bounded, monotonic, and zero iff no self-reference. Three blind spot theorems establish provable limits.

\bibliographystyle{plainnat}
\bibliography{references}

\end{document}
